\section{The Metaphysics of Sex and the One}

\paragraph{Interim Note on Being One with Everything}


\begin{quotex}
Pantheism does not admit the existence of individual beings.

Pantheism offers only the perspective of letting oneself be lulled by the undulation of the ocean of deified nature. \flright{\textsc{Valentin Tomberg}}

\end{quotex}

There is a photo of \textbf{Giulio Cogni} in \textbf{Julius Evola}'s \emph{Sintesi di dottrina della razza} with the following caption:

\begin{quotex}
An Italian Nordic-Aryan type of the body with an Amazonian race of the spirit. Actually, it is a matter of a person, in whose theories an external energetic and almost Promethean aspect is mixed with a Demetrian-Lunar pantheistic vision of the world with a special recognition given to the feminine element. 

\end{quotex}
Giulio Cogni was a philosopher, music critic (a specialist on Wagner), and an author of books on Eastern spirituality. After the German-Italian alliance in the 1930s, there was a scramble among Italian intellectuals to outline a Fascist understanding of race, in opposition to the German's. The consensus was a more spiritual understanding of race. Thus Julius Evola published several works on the topic. His young friend, Massimo Scaligero, published a book on race based on some ideas from Anthroposophy. Giulio Cogni, who was a follower of the philosopher Giovanni Gentile, tried to reconcile Alfred Rosenberg with the situation in Italy. From the snippets I've found, Cogni denied the notion of racial superiority and, not unlike Evola, claimed that ``Aryan" referred to a spiritual quality rather than a biological quality.

After the defeat of WWII, Evola more or less abandoned the topic. Scaligero and Cogni, on the other hand, managed to rebrand themselves: Scaligero as an Anthroposophist and Cogni as a Vedantist. As an indication of their success, some of their articles have been republished in the American \emph{Yoga Journal}, with no mention of their Fascist past.

After the war, Evola's and Cogni's life paths diverged. I did find a rather long review by Evola of Cogni's book \emph{Io sono te: sesso e oblazione} (``I am you: sex and offering"), which highlights the differences between the two thinkers and would explain Evola's evaluation of Cogni.

In that book Cogni proposes a pantheistic worldview based, as he claims, on the Vedanta: ``we are all one". That logically leads to an undifferentiated view of the human situation, which includes promiscuous sex, self-sacrifice, and androgyny. Since Cogni articulates in an intelligent and sophisticated way many of the commonly held views in the modern world, the review is an apt essay for discussion. 

\paragraph{Being}
Cogni bases his ideas on the notion of the Vedantic One, as he understands it to indicate the overcoming of all distinctions. Quoting Hegel, Evola describes this view as ``the night in which all the cows are black". I noticed this years ago when I participated in a \textbf{Nisargadatta} discussion group. The consensus view in that group was that, as an enlightened being, he more or less disappeared into the ``One". Nevertheless, as I pointed out, he spoke, responded to his name, etc., but that had little impact on the discussion.

In that view, there is only undifferentiated Being, so the beings we experience are illusions. ``Whose" illusions they are, or why those ``illusions" take such a specific form, or why do different people experience a similar world, was never made clear if only the One exists.

\textbf{Jacques Maritain}, in \emph{The Degrees of Knowledge} provides a more satisfying answer. When I know Peter, I judge him to be a man, but even more fundamentally, I have an intuition of Peter as a being. This intuition is found everywhere, it ``saturates all things". Maritain explains our experience of Being.

\begin{quotex}
Being is a primordial and general conceptual object which is at once and from the beginning essentially diverse in the diverse subjects in which the mind discovers it. What is primarily known, and in which every object of thought is resolved for the intellect, is being. But nothing can be added to it extrinsically to differentiate it, for all its differentiations issue from its own depths. 

\end{quotex}
In other words, Being is known through beings. \emph{Pace} Cogni, Being is already highly differentiated. In the Western tradition, this is expressed as the doctrine of the consubstantiality of the Logos with Being. Evola expresses the same notion in a slightly different way:

\begin{quotex}
The One dominates a well-articulated order of differences (a cosmos, in the Hellenic sense). 

\end{quotex}
So, ``Illusion" has a much different connotation. For Cogni, the illusion is the separateness and distinction of beings, which are ultimately part of Brahman. However, illusion is actually much more subtle, and hence, more difficult to eradicate. The illusion arises from our failure to grasp the meaning of the Whole and how the parts are interconnected. This is due to the weakness of the intellect and erratic emotional reactions.

\paragraph{Descending and Ascending Movements}
Evola claims that Cogni fails to distinguish between the subconscious and the superconscious, with the following consequence:

\begin{quotex}
Cogni has no sense of the fact that just as an ``integrative ascending self-transcendence" exists, so a divisive descending self-transcendence also exists for the true personality. That is, there are possible openings of the I both toward the higher as well as toward the lower, which also means toward ``nature", toward the unconscious, toward the vital formless bottom. 

\end{quotex}
In an early essay, \emph{The Holy Spirit \& the Sophia}, Tomberg describes that downward movement in these steps.

First the Logos descends to the conscious I. If the I is oriented to the Logos above, it is ascendant. This supersensory ``Temple of Wisdom" or ``House of Sophia" is the knowledge of the Whole (but not of every individual fact).

However, that movement can continue downward in these stages, each with its own type of intelligence:

\begin{itemize}
\item \textbf{Hunger}. The stream of the Logos is brought down to the level of digestive activity. This leads to a logic in which the stomach and its interests determine logic and conclusions. This is the level of Marx. 
\item \textbf{Sex}. The stream is brought down to the instinctive life. This is the level of Freud and psychoanalysis. 
\item \textbf{Nature}. The stream descends to the subterranean spheres. This logic may seem quite convincing. At its lowest point, the individual no longer thinks independently, but thoughts arise through him. 
\end{itemize}
\paragraph{Pantheism}
Tomberg distinguished between theism, pantheism, and naturalism, which are valid on different planes:

\begin{itemize}
\item \textbf{Naturalism}: only what is experienced through the senses is real 
\item \textbf{Pantheism}: thought alone can grasp the world. Existence is illusory, only the Absolute, One, Brahman, or God is real 
\item \textbf{Theism}: Transcends rational thought. Existence not strictly logical, i.e., cannot be derived from thought. 
\end{itemize}
Evola offers something similar in this comment:

\begin{quotex}
So Cogni denies that that polarity is an essential requirement of eros, which he believes concerns only the naturalistic plane, like electrical and similar phenomena. That means that for him all the documentation that we have gathered, in an entire chapter, from the most varied cultural areas, regarding the ``metaphysical dyad", might as well not exist because it contradicts his promiscuous pantheism. 

\end{quotex}
Thus, on the naturalistic plane, any sexual experience is worthwhile in its own right if it is observed in natures. Neither electricity nor eros depends on a anything higher. For pantheism, the male-female polarity is just a form of thought with no ultimate reality. It is no more than one particular expression of sexuality.

Not to forget that the young Evola's first book was a translation and commentary of the \emph{Tao Te Ching}, he understands that, on the metaphysical level, the Yin-Yang, male-female, dyad is fundamental. Hence, sexual expression on the physical plane needs to mirror that.

\paragraph{Postscript on Sexual Intelligence}
The flip side of Evola's critique is that Cogni has empirical experience and contemporary mores on his side. As we saw above, the sex instinct has its own logic and intelligence. Those who have ``transcended" the self in the descending direction will be baffled by appeals to sexual restraint based on higher principles.

The occultist \textbf{Alice Bailey }seemed to have had an intuition of this, making predictions that have largely come to pass. In \emph{Esoteric Healing}, she wrote:

\begin{quotex}
Inter-marriage between nations and races, the fusion of bloods for hundreds of years—due to migration, travel, education and mental unity—has led to there being no really pure racial types today. This is far more certainly the case than the most enlightened think, if the long, long history of mankind is considered. 

\end{quotex}
The European wars of last century, not to mention more recent turmoil, have certainly exacerbated those four factors. Eros, or the sexual drive, is potent, even to the point of overcoming barriers:

\begin{quotex}
Sexual intercourse knows no impenetrable barriers … The urge to mate becomes peculiarly strong when men are removed from their familiar settings and experience the novelty of complete loneliness, when the normal inhibitions and customs imposed by family relationships and national standards are removed, when danger of death is constantly faced and the larger value submerges the lesser values and the usual conventional attitudes 

\end{quotex}
\paragraph{Appendix on the 120 Days of Sodom}
Curiously, Evola accuses Cogni of tending towards the sexual excesses of the \textbf{Marquis de Sade}. This may not be so far-fetched, given the tendency of some on social media to associate the ``right" with bondage, domination, and so on. They are not to be taken seriously.

Nevertheless, the Italian communist film director \textbf{Pier Paulo Pasolini} created a work based on de Sade's book, but situated it in the Salo Republic. Perhaps he derived that association from Cogni, but certainly not from Evola who rejected such practices. Paradoxically, it was Pasolini himself who advocated sexual experimentation. He was murdered by an underage hitchhiker who objected to Pasolini's proposition to sodomize him with a broomstick.

\paragraph{Translation}

\begin{quotex}
Evola's review of Giulio Cogni's book \emph{Io sono te: sesso e oblazione}
(You and I: sex and offering). 

\end{quotex}
In general, we do not believe that discussions have a meaning and a seriousness unless it essentially aspires to a
clarification in the premises of a common base. If a writer is able to recognize the fundamental premises of his own
thought (whether or not tied to a ``personal equation") and sees his fundament difference in respect to those of another
writer, the only serious thing is that he follows his own way without trying to insert himself in an intellectual world
foreign to him. Unfortunately that happens rarely, due to the lack of such a preliminary self-analysis. That does not
limit us to an inescapable criticism of other people's views (which is of course approvable and
fruitful), but is put on a confused argument: precisely because of the heterogeneity of basic conceptions. That
demonstrates fundamentally the activity of subrational explanations.

For this reason, we would not have considered that a book titled \emph{Io sono te: sesso e oblazione}\footnote{\textit{You and I: Sex and Offering.}} was by Giulio Cogni. The author has certainly believed he included in it an essay on our work,
the \emph{Metaphysics of Sex}. That would not be as important as the fact that in the circumstances grave confusions
and distortions weighed on a domain that goes beyond the ideas that we've espoused in various
places. Hence, the opportunity of a development that, apart every polemic need, intends to highlight some ensemble of
ideas which may be of interest for the reader.

The views of \emph{I am you} reproduce, in its essentials, that which Cogni had already expressed in another
book, \emph{Saggio sull'amore} [Essay on Love], published in 1922. At that time, someone wisely
said that Cogni, at that time a Gentilian through the skin, had translated the Gentilian theory of the ``spirit as pure
act" into a more sapid ``theory of the spirit as impure act", since he recognized in the sexual union a principle
concrete form of the identification of the subject with the object postulated by Gentile's
actualism. Moreover, Cogni formulated a phagic or anthropophagic theory: love, would mean eating oneself, devouring
oneself. A similar thesis in his new book: where ``the famous hunger — sex, hunger-love, equation"
is established. Previously, Cogni presented the situation in rather masochistic terms: man eaten by the woman in whom
he was absorbed and lost himself as individuality. In the more recent exposition erotic anthropophagy seems to be
conceived as reciprocal, turning out however to be little imaginable because, ironically, he would want to think that
at the end of the lovers, there remained only two mouths, each being ingested and consumed wholly by the other.

If it all ended here, with phagia, it would be said the Cogni was inspired only by the crudest aspect of sexuality: the
``hunger" of bodies, simple lustful longing. But suddenly he proceeds to appositions that are in complete contrast with
what, as analogy, might suggest ``eating" and hunger. In fact, Cogni always returns again to self-denial, to sacrificial
offering, the abandonment of oneself that there would be in eroticism and in sexual union. Carnal pleasure would be
``complete renunciation of oneself in order to make oneself the other."

He moves directly to a type of mysticism: ``the sexual act is humility of annihilation and sacrifice of oneself to the
universal life that is seen in the body of the beloved". ``Love is phagic because only with the gift of the body and
with devouring and making oneself devoured one realizes the more powerful symbol of Unity: every alleged individual
separation is removed." The ultimate goal is ``the immersion of the cosmic One without a second." All this seems to us
pure digression, even with a slight paranoid tint.

First of all, about the incongruence between the points of view, it must be pointed out that in hunger and in phagia
presented as the keys of sex, there is nothing of this sacrificial orientation, of the abandonment of oneself, and of
``sweet" identification (sweet is a word that often recurs in Cogni even in connection to sadomasochistic situations;
that would point, to malign him, at the strong preference, usually female, for products of the confectionary industry
instead of the prevalent manly taste for spicy and peppered food). indeed, in hunger the need operates by suffering
from a deprivation, while eating leads only its own conservation and satiety: exactly the opposite of the abandonment
of the sacrificial gift of oneself. Nothing sweet, then, if hunger is absolute, devouring. And since Cogni even makes
almost a dogma out of anthropophagy in its proper sense, he does not know that in it, at most, the situation is not
different: it is multiply attested that if the savage feeds himself with the flesh of others, he does not do it for a
``sweet" identification, but only and darkly because he believed he absorbs the other's strengths
to this own advantage. As for eating the flesh of sacrificial victims, it is a fantasy to say that in it the tendency
acts to submerge oneself in the cosmic One. In general, here everything comes down to totemic participation (the victim
incarnates the totem), then a rather restricted order, shadowed by sorcery and demonism typical of totemism in general.
Therefore, these are very fragile and inconsistent bases for the advocated theses referring to sex. And the Eucharistic
symbol, if one does not want to totally contaminate it by discovering in it rather suspect roots, is reduced to a mere
allegory.

As far as we are concerned, the fundamental intention of the book \emph{Metaphysics of Sex} was to highlight the
existence of a possible transcendent dimension of sex. We tried to lead the transcendental meaning of eros (in the
almost Kantian sense) back to a dark, unconscious impulse to restore an original wholeness, which we referred as the
most noted mythical formulation in the West, to the Platonic theory of the androgyne. Besides, we found that phenomena
of transcendence could intervene in the erotic-sexual experience of a momentary traumatic removal of the common
conditionality of the individual consciousness, we showed that this was the presupposition of practices of some
environments, especially the Oriental, in which it was done in a magical, initiatic, or evocatory use of sex. But all
that is quite far from Cogni's mystical-phagic digressions, and every confusion in this regard is
deplorable.

First of all, we should not generalize by attributing the strong value of transcendence to that which is typical of
almost all sexual unions of human beings. What appears in the metaphysical or transcendental cannot be made of value in
the phenomenological. In a phenomenological examination fall all the idealizing and mystical fixations, with the very
sweet sacrifice and the sacrificial surpassing of oneself in the flesh of the other which Cogni always comes back to.
Factually, in most sexual intercourse, one partner seeks only his \emph{own} pleasure, making the other the means, so
that the situation is not very dissimilar, \emph{sit venia verbis,} from a ``masturbation duet". Therefore nothing
overcomes the individual bond. In the second place, existentially, often, and today more than ever, sex is of value to
the individual only as a self-confirmation, only to assert a ``need" (\emph{Geltungstrieb} as Adler would say), or to
look for an illusory, turbid substitute for a true sense of existence, of which he is deprived: then, again, no exit
from the closed circle of the individual. Finally, as we said and also highlighted in our book, if sporadic phenomena
of ``transcendence" in the profane experience of sex can sometimes be verified, they often are not \emph{experienced} as
such. They are realized in apical forms to the trauma of the organism which represent, for the most part, a solution of
continuity of the consciousness from which one returns emptied, instead of having had the ``shining experience of the
One". That, in general, is rarely the case in profane, carnal, or romantic love. It is especially pertinent to the
magical and initiatic use of sex, the use that, inter alia, involves a special conduct of sexual intercourse and about
which one thing is certain: intensive states of a special destructive intoxication intervene (in an almost ontological,
not a moral, meaning), which excludes phagia, the abandonment of the other, the sacrifice of oneself, and all the sweet
and pantheistic affectations so dear to Cogni; and they never failed to emphasize the dangerous character of that
practice, for anything idyllic, romantic, and idealizing.

Cogni has also noticed the relation, noticeable in multiple forms, between eros and death, between the divinity of sex
and the destructive divinity and death, that we did, but without understanding it in its true sense. It is significant
that, among other things, the secret Hindu orgiastic rites that aimed at the aforementioned experience of
transcendence, was celebrated in the sign of goddesses like Durga and Kali, not in their maternal aspect but in the
destructive aspect. Sekhmet, the Egyptian goddess of love, is also the goddess of destruction and war (her leonine head
refers to a beast whose manners are not at all the sweetest). Something analogous applies to the ancient goddesses of
the Mediterranean area, starting with Ishtar, also the goddess of the orgy.

The following point is related to that. We indicated the magnetic base of every eros and every intense sexual
experience. It is the development of such a base that serves as the basis for those experiences. But it is due
precisely to the polarity of the male and the female as ontological principles, something always acknowledged. So Cogni
denies that that polarity is an essential requirement of eros, which he believes concerns only the naturalistic plane,
like electrical and similar phenomena. That means that for him all the documentation that we have gathered, in an
entire chapter, from the most varied cultural areas, regarding the ``metaphysical dyad", might as well not exist because
it contradicts his promiscuous pantheism.

Without dwelling on such a, perhaps a little too specialized, topic, let's indicate in general the
horizons between which the erotology proposed by Cogni moves. The theory of eros as pantheistic identification has
found nourishment, in Cogni, in his references to India and to Hindu philosophy of the Vedanta. It appears in a clear
way that Cogni has seen of India only what, given his temperament, he was capable of, the India that would be immersed
in the ``dream of the One", the ``hot all comprehensive all justifying sweetness, royally tolerating, loving and
accepting of the people and the land of India". \emph{Mother India}, discovered by certain useless humanitarian
American authoresses like women, with Gandhiism, non-violence, the alleged climate of ``loving equality" for the meaning
of the One above every illusory difference, would be the true India. But that is an image in part one-sided, in part
absurd.

First of all, Cogni seems to overlook the small defect of the beauty present in the ``sweet tolerance", exhibited by the
recent massacres between the Hindu and Islamic inhabitants of India, which Gandhi made out as another pleasant event.
He then ignores the fact that, if there was a social regime which ruled in the most rigid way the principle of
difference for entire millennia, this was the Hindu caste system. To the alleged India all love, abandonment and pardon
we contrast the India of the great epochs and of the Bhagavad Gita, a traditional text having the same popularity which
the Bible enjoys among the Westerners. It attributed the characteristic of a passion, destructive in its transcendence,
to the highest form of the appearance of the divine, while also drawing from it a spiritual and metaphysical
justification of the duty of the warrior to fight and kill, sparing neither friends nor relatives who were found on the
enemy front. And everyone knows that the Hindu Trimurti, much closer to the Hindu population than the abstractions of
Vedantic speculation, considers a divine function in Shiva to whom destruction is proper.

But now it is necessary to repeat myself that Cogni is visibly affected by spiritual scotomas which prevent him from
seeing whatever does not support his inclinations. So Buddhism interests him only in its late and popular exoteric
forms of a religion, with ``love for all creatures", Amitabha the god of love, etc., in contrast to the serious
individual ascetic techniques of the Buddhist doctrine of ``awakening" of the authentic texts of the original Buddhist
canon, that we showed with strict reference to the texts in a book that Cogni says he was familiar with. From that
canon it becomes obvsious that, among other things, if love and compassion figure (however with an instrumental
function) as preliminary stages in the sequence of the four phases of the highest Buddhist contemplation,
of \emph{dhyana}, they are left behind, since the summit is constituted by a state of sovereign, disincarnated
impassibility and imperturbability that, pleasing or not to Cogni, has something of the Olympic quality and nothing of
a week humanitarianism.

In fact, our author does not stand at the peaks but in the trash heaps of India. The current devotion has had a part in
India, but in the lower popular strata, not unrelated to a pre-Indo-European substrate of the country. Only relatively
recently has it corresponded with a philosophic system, that of Ramanuja. Earlier, it was considered a ``path of
devotion" and love, \emph{bhakti-marga}, but it certainly did not stand at the first level, the dignity of a spiritual
``regal path", \emph{raja-marga}, which was instead attributed to the path of
knowledge", \emph{jnana-marga} and \emph{jnana} yoga.

That character was mainly attributed to the Vedanta, which Cogni is enthusiastic about, seeing in it however only the
theory of absolute Identity, of non-duality, of the One-All, of ``thou art that", the theory that serves him as the
basis for his ideas of eros that embraces and reunites everything.

Well, here we can first of all say that the Vedanta in the primitivistically pantheistic key does not at all exhaust the
Hindu spiritual world. We can point out, especially, that a radical monism was not attested in its origins, in the
Vedas, which present us with a clearly articulated \emph{pantheon}. In the second place, India has known great
speculative systems, like the Samkhya, which has instead emphasized a primordial duality, that of Purusha and Prakriti,
and like metaphysical tantrism, which argued against the ``illusionistic" Vedanta (the world is \emph{maya}) and along
with the school of Kashmir formulated a well differentiated cosmological doctrine.

We will not dwell on these factual data about India. The essential point is that Cogni exchanges the metaphysical One
with the pantheistic One, with that One which, according to the expression Hegel used in regard to the philosophy of
identity of the later Schelling, is ``the night in which all the cows are black". It is not about the One that dominates
a well-articulated order of differences (a cosmos, in the Hellenic sense) but rather of a promiscuous ``naturalistic"
unity to be related, rather, to ``Life". This is Cogni's spiritual horizon.

From that confusion, much more serious confusions arise in the practical realm. Cogni has no sense of the fact that just
as an ``integrative ascending self-transcendence" exists, so a divisive descending self-transcendence also exists for
the true personality. That is, there are possible openings of the I both toward the higher as well as toward the lower,
which also means toward ``nature", toward the unconscious, toward the vital formless bottom. Only the former correspond
to the high ascesis, to initiation, to authentic yoga. The ancient wisdom already distinguished and counterpoised the
higher waters from the lower waters, the first illuminating, the latter intoxicating and divisive, and this basic
doctrine, taken up again by thinkers around the period of the Renaissance, was opportunely recalled by one of the few
people today truly qualified in this field, by Rene Guenon, to alert us about the danger and deviations of a certain
contemporary spirituality.

To return to the field of eros, Cogni recalled in passing the ambivalence that is present, from the spiritual point of
view, in sexual experience. If the woman was seen as a danger, if one could say \emph{foemina mors animae} in Latin so
as to recommend continence, that is not attributable to a moralistic attitude, to the ``theological hatred for sex" that
Vilfredo Pareto spoke about or to the ``sexophobia" on which L De Marchi insists. They also had in view the possibility
that the experience of sex, when it was not about curbing mere mortals and moralizing, but for those who had
supernatural aspirations, that it could even lead to the negative direction of a ``descending self-transcendence". And
if we examine the use of sex predominant in the most recent generations, we find a reflection of that even on a quite
profane plane: other than a ``sweet" sacrifice of oneself, a ``carnal offering", that redeems and leads to the One, but
couplings used on the same life of drugs (to be precise: of the current profane pandemic use of drugs) to draw from the
extreme sensations of the orgasm the illusory confirmation of the sense of oneself (the exact opposite of the direction
toward the higher).

When the One is ``the night in which all the cows are black", every difference is challenged and discouraged, and
promiscuity in the name of that One becomes a norm even in the forms more repugnant for every well born person. In that
regard, Cogni is explicit and demonstrates, if nothing else, the courage of coherence. He asserts, for example that
``every doctrine that starts from the absolute and not relative point of view, from inferiority or hierarchal
superiority, is erroneous at its base, if it is true that the One is all and identical to Brahman". Note that these
words are said in regard to the difference between the species, between men and animal species, for example. We are not
talking about how it relates to the human environment. Cogni will certainly not object, to the greater glory of the
Vedantic One, if a young Nordic girl beds a Zulu or an Australian aborigine whose morphological and mental level
corresponds to the stone age. He is certainly an egalitarian to the bitter end, a fanatic integrationist (he rushed to
make an act of contrition, keeping to date in relation to a past error because in the fascist period he was a
racialist, albeit of a dubious racial theory), an admirer of ``unisex" and the ``third sex" and so on. But the last straw
is in the area of carnality. As a more audacious form of identification, phagist or not, in the name of the One, he
actually admits not only homosexuality and pederasty, but even sex with animals, sodomy with women, and so on. His
theory explains how ``many people have so much attraction for such relationships commonly judged against nature". ``Only
by accepting in principle the ordinarily most repugnant areas of the other (for sexual use), is one certain of having
reached absolute identity." At this rate, we believe that Cogni would, in the same way, value even coprophagia (the
eating of feces as a form of eroticism) and would sanctify the disgusting events where even coprophagia figures
abundantly as the decisive erotica along with the other horrors, described in the \emph{120 Days of Sodom} by
the \textbf{Marquis de Sade}.

Naturally, in denouncing similar aberrations of Cogni, we do not appeal to any conformist moralism but to what is called
normal in a higher, not social, sense. For example, pederasty at the most can be tolerated when if arises from special
constitutional situations of imperfect sexualization but must be stigmatized as a vice, deviation, and perversion in
all other cases, denying also that in this, as in all other case of sexual psychopathy, the necessary objective
requirements from a metaphysics of sex are present for the actual experiences of a deconditionalizing intensity. But
there is no hope that Cogni has such an understanding for similar things.

Finally, we need to point out another deviation in Cogni, unifying promiscuous sexuality with pantheism. Since, as we
have seen, his reference point is not metaphysical reality but rather the promiscuous bottom of ``Life", skirting the
unconscious and the subconscious; in his more recent writings, Cogni sympathizes openly with psychoanalysis and
metapsychics. He goes so far as to say that parapsychology ``remains still the great hope of the future". He poses
himself in a duet of mutual admiration with Emilieo Servadio who had ``opened the way", ``knowing India in depth and
every type of initiatic thought and psychic depth". That makes us smile. If Servadio had ever really had some idea of
intiatic matters and authentic wisdom that was when, before the war, he was vividly interested in the publications
edited by the ``Ur Group" which I directed. After the war, at a snail's pace, he more or less set
all that aside and he immersed himself in psychoanalysis. Moreover, that brought on a lucrative professional program
and seeking to put himself everywhere before the public, he associated psychoanalysis with parapsychology, in place of
initiatic knowledge and wisdom traditions, giving himself however to probe not the psychical depths but rather the
slums of the psyche. The connection with psychanalysis and metapsychics took place in Cogni without difficulty through
the fact that his One is easily relatable to the ``deep unconscious that is one in the entire universe", by which
telepathic phenomena and metapsychics in general can be explained, while it is the field proper to psychoanalysis.

In fact, in these new disciplines the unconscious becomes a sack in which things of every type are placed. One does not
think at all about such an elementary and basic distinction, like that between the subconscious (or unconscious) and
the superconscious, in part for the simple reason that psychoanalysis and metapsychics have no idea of the latter, and
then because in their experimental field, for obvious reasons, it is rather difficult to even imagine. Form this also
come aberrant assimilations, like those of C G Jung, in which the homologous figures perceived by psychopaths or
oneiric visions of symbols and mythical structures of the initiatic and religious field, reducing everything to the
emergence of the archetypes of the collective unconscious. Now, it still holds that, apart from psychoanalysis with its
murky world, all modern parapsychology embraces only the offal of the extranormal and is outside all that might have an
authentically spiritual value. It is only a question of the slums that can impress only the naïve. But the reason for
Cogni's interest in such things appears evident given what is said: it is a question of true
elective affinity. At this level is where he locates which he called the ``great hope". Signs of the times.

In conclusion, from the present considerations it turns out that the mentality, the personal equation, the elective
affinities, the theoretic referents of Cogni have no point of contact with a spiritual world that we presume is not
personal. In cases of this type, as was said at the beginning, discussions have little meaning. However, these notes
are less about ``arguing" than taking the occasion to make clear certain questions of undoubtable relevance for those
readers who might have an interest in the themes treated.

\flrightit{Posted on 2015-10-08 by Aeneas }


\begin{center}* * *\end{center}

\begin{footnotesize}\begin{sffamily}



\texttt{Aleksandar Gueric on 2015-10-15 at 01:40 said: }

Alice Bailey certainly knew all about the united world and homogeneous raceless population agenda due to the connection between Theosophy and various groups of shady origin, including Fabian society, who, in the spirit of Bavarian Illuminaty, contined the work of subtle social engineering towards the Age of Aquarius i.e. Lucifer. We are all witnessing a progress of this plan and it's acceleration. I wouldn't be surprised if they had connection with the Frankfurt school as well.

Speaking of ascending and descending in regard of transcendence of human experience I am reminded of Terrence McKenna's ``Archaic Revival", a concept of complete cultural and individual dissolution of boundaries and return to, what he called, protohistoric paradisical existence of pastoral, orgiastic and classless mushroom worshipping potheads. McKenna was a feminist, profound one some may say, but still a feminist who never missed a chance to blame evil patriarchy and it's fascist spirit for all the wrongs in modern world. Well, his idea of the return into a sacred womb of matriarchy included different kinds of human descendance into the subconscious and nano existence within the cyber world. On the other hand, he considered ascending spiral ``techocratic and fascist". When we research his origins we will find the same shady connections. Since I don't intend to turn this comment into a conspiracy rant I will stop here. But there is definitely a global agenda regarding descending transcendence and it's popularization.

Italian racialists had every right to criticize National Socialists race theory which ignored spiritual dimensions and focused strictly on blood alchemy. How many Aryan looking nordics are out there, with the completely atrophied spirit and essence? Could we call them Aryans? Certainly not, for Arya is a spiritual nobility as well. Arya is a title not a race, although it's reserved only for certain races (indo-european stock).


\hfill

\texttt{Cologero on 2015-10-16 at 07:42 said: }

Cogni's views are more in tune with the signs of the times. His views on sexuality, psychology as a spiritual path, pantheism, paranormal experience (e.g., the widespread interest in ``channeling", astral travel, shamanism, the ``law of attaction") have become mainstream since Cogni's book was published. On the other hand, Evola's insistence on the transcendent ``metaphysical dyad" is losing steam (and is not aided by Evola's past association with race and fascism). Even one of the last holdouts for the dyad, the Roman Church, is wavering, if the discussions emanating from the Synod on the Family are any indication.

The Italians started destroying everything associated with Mussolini as the Allies advanced toward Rome. Thus, it is impossible for an independent scholar to find most of the books from that era. However, I did find this article, \textit{The Italian Roots of Racialism}\footnote{\url{http://escholarship.org/uc/item/85r1607t\#page-1}}, which has a discussion on Cogni. The writers discussed may have well as been space aliens, since their viewpoints are so far from anything being expressed today.


\hfill

\texttt{Sparrow on 2015-10-18 at 06:39 said: }

Since he was a Vedantist, one would think that Cogni read the Bhagavad Gita and realized that the ``Oneness," is not a call to promiscuous dissolution rather than a virile call to action and obeying one's dharma.


\hfill

\texttt{Cologero on 2015-10-18 at 15:37 said: }

Yes, Sparrow, that is precisely the point Evola made against him. Since I haven't been able to find Cogni's text anywhere, it is impossible to make a complete judgment. Nevertheless, he seems to be among the ``neo-vedantists" who are prominent today in the West. (A google search will provide many examples.)

Since Cogni was heavily influenced by psychoanalysis, I suspect he had in mind Freud's notion of the ``oceanic feeling" in which there is a return to the infantile consciousness of being ``one with the world".


\hfill

\texttt{SB on 2015-10-18 at 15:55 said: }

Is Cogni very well known among spiritual and philosophically inclined people today in Italy? And if so, is his reputation more ``mainstream" than Evola `s ?


\hfill

\texttt{SB on 2015-10-18 at 15:59 said: }

Re. Pasolini, if I may offer a comment re. his association of Sadism with Fascism, I believe it may simply be a result of New Left association of authoritarianism with sexual dominance, as part of their general worldview basing everything on sexuality (derived ultimately from Freud but more to the point from Wilhelm Reich, Herbert Marcuse etc )


\hfill

\texttt{Cologero on 2015-10-18 at 16:46 said: }

@SB, as I mentioned, Cogni's books are unavailable in print or electronic editions. Hence, I can't imagine that he has much of a reputation at all anymore. What I did say is that his worldview appears to be more congenial to people today. For example, he is much more open to, let's say, the ``varieties of sexual experience" than Evola was. Evola was more attached to the boy-girl thing, not for any moral or religious reasons, but rather because he believe it reflects the polarity of the male-female, or yin-yang, princples, the overcoming of which results in the ``androgyne".

As for your other point, Pasonlini, Marcuse, etc., seem passe now. Such an association was shattered by Michel Foucault's life and work. He combined gay S\&M with libertarian economics, sort of the precursor of the ``socially liberal but economically conservative" type, not uncommon today. But he went further, challenging power structures that would marginalize certain groups, perhaps those who could be considered the ``fifth caste" that the Tantras refer to.


\hfill

\texttt{SB on 2015-10-18 at 18:23 said: }

If I'm not taking the discussion off topic, what's the fifth caste as I'm ignorant of the subject? The pariah/untouchable? 

In response to your mentioning the New Left as being passé, my observations would be that society has hit a key turning point with the transition to postmodernism; the rationalist mentality that has been dominant since the XVII-XVIIIth centuries, the search for a theory which explains everything, has been judged a failure and now rather than proposing new systems or axioms for all to follow as they did in the past, philosophers are questioning every existing remaining ``given" rule of society with the aim of abolishing it, to the extent where language even no longer has a definite meaning (Derrida and Deconstructionism). It's an interesting thought process to see where it's leading to; one could well say the world is heading to the submersion into the undiferrentiated all . The rise of the environmentalism with their quite pantheistic beliefs (Gaia hypothesis) is another trend in this direction.


\hfill

\texttt{Aleksandar Gueric on 2015-10-18 at 19:22 said: }

"If I'm not taking the discussion off topic, what's the fifth caste as I'm ignorant of the subject? The pariah/untouchable?"

You probably think of dalits. European gypsies descend from them, which is proven in genetic studies.


\hfill


\end{sffamily}\end{footnotesize}
